\clearpage
\subsection{一个辅助链复形}\label{sec:17e}
另取正合 Lagrangian 子流形 $Q$。在此前选择之外，为子流形对 $(Q,L_k)$，$k=0,1,2$，固定 Floer 数据，并固定定义三角形乘积
\[
\begin{aligned}
\mu^2 &:CF^{\mathrm{pr}}(L_0,L_1)\otimes CF^{\mathrm{pr}}(Q,L_0)\longrightarrow CF^{\mathrm{pr}}(Q,L_1),\\
\mu^2 &:CF^{\mathrm{pr}}(L_1,L_2)\otimes CF^{\mathrm{pr}}(Q,L_1)\longrightarrow CF^{\mathrm{pr}}(Q,L_2)
\end{aligned}
\]
所需的扰动数据，以及与这些数据一致的四标点圆盘万有族上的数据，后者给出
\[
\mu^3:CF^{\mathrm{pr}}(L_1,L_2)\otimes CF^{\mathrm{pr}}(L_0,L_1)\otimes CF^{\mathrm{pr}}(Q,L_0)
\longrightarrow CF^{\mathrm{pr}}(Q,L_2).
\]
略加改变上述想法，即可定义映射
\begin{equation}\label{eq:17.15}
\boxed{\begin{gathered}
d:CF^{\mathrm{pr}}(Q,L_1)\longrightarrow CF^{\mathrm{pr}}(Q,L_2),\\
l:CF^{\mathrm{pr}}(L_0,L_1)\otimes CF^{\mathrm{pr}}(Q,L_0)\longrightarrow CF^{\mathrm{pr}}(Q,L_2),\\
\mu^1(d(\cdot))+d(\mu^1(\cdot))=0,\\
d(\mu^2(\cdot,\cdot))+\mu^1(l(\cdot,\cdot))
+l(\mu^1(\cdot),\cdot)+l(\cdot,\mu^1(\cdot))=0.
\end{gathered}}
\end{equation}
为以更紧凑的形式编码这些关系，定义
\begin{equation}\label{eq:17.16}
\begin{gathered}
D=CF^{\mathrm{pr}}(L_0,L_1)\otimes CF^{\mathrm{pr}}(Q,L_0)
\oplus CF^{\mathrm{pr}}(Q,L_1)\oplus CF^{\mathrm{pr}}(Q,L_2),\\
\partial=\begin{pmatrix}
\mu^1\otimes\Id+\Id\otimes\mu^1&0&0\\
\mu^2&\mu^1&0\\ l&d&\mu^1
\end{pmatrix}\in\End(D).
\end{gathered}
\end{equation}
于是\eref{17.15}等价于 $\partial^2=0$。映射 $d,l$ 实际上就是构造 Floer 上同调长正合列所用的映射，见\bcite{131}第 3.3、3.4 节，\znote{247}{\eref{17.16} 是长正合列的“链级包装”：$\partial$ 的第一行是 $\mu^1$ 在张量积上的 Leibniz 作用，第二行是 $\mu^2$，第三行是 $l$ 与 $d$；$\partial^2=0$ 与\eref{17.15}的两式逐一对应。这种写法的好处是把“$\mu^1$ 加更高阶修正”压成单个微分，从而可以整体做线性代数：下文与代数一侧的锥比较时，正需要把 $(d,l)$ 放进矩阵的第三行，而不是逐条处理关系式。[来源：原书 §(17e)，印刷页 247–248。]}
\SourcePage{256}
而\eref{17.16}在其证明中起关键作用，见\bcite{131}引理 2.32。不过，由于这里与\bcite{131}之间有记号及少许技术差异，仍值得简要回顾定义。

对于 $d$，从 $\pi:E\to S$ 出发，不仅去掉 $\zeta^{d,-}=*$，还去掉另一个边界点 $\zeta^{d,+}$；例如从 $*$ 沿正方向稍移一点得到的点，这样做的好处是纤维化在该点附近仍局部平凡。结果为 Lefschetz 纤维化 $E^d\to S^d\cong Z$，有一个入向条带状端和一个出向条带状端。相应边界条件 $F^d$ 如下定义：在“短”的边界分支上恒定等于 $Q$；在“长”的分支上，$\epsilon_{\zeta^{d,+}}(s,1)$，$s\geq0$，上方纤维中为 $L_1$，$\epsilon_{\zeta^{d,-}}(s,1)$，$s\leq0$，上方则为 $L_2$，利用平行移动及同痕 $\tau_V(L_1)\simeq L_2$ 在两者间插值。\sr{fig:17.4}{图~17.4}给出该纤维化的两个图景，只在底空间的画法上不同。赋予相对扰动数据，并定义 $d=C\Phi_{E^d/S^d}$ 为所得不变量。
\SourceFigure{figure-17-4}{0.52546097}{17.4}

$l$ 的构造与 $k$ 的构造十分相似。取三标点圆盘上的平凡纤维化，带边界条件 $(Q,L_0,L_1)$，将其入向端与 $E^d\to S^d$ 的出向端粘合，得到纤维化 $E^{l,1}\to S^{l,1}$，带边界条件 $F^{l,1}$ 及相对扰动数据。另一方面，可取三标点圆盘上的另一个平凡纤维化，这次前两个边界分支上的边界条件分别恒定为 $Q,L_0$，第三个分支上则实现从 $\tau_V(L_1)$ 到 $L_2$ 的同痕。在 $L_0$ 分支上标记一个边界点，再在 $*$ 处与带标准边界条件 $F$ 的纤维化 $E\to S$ 作边界连通和，得到另一个纤维化 $E^{l,0}\to S^{l,0}$；如前，\stref{笔记}{17.8}适用。存在明显的族 $E^{l,r}\to S^{l,r}$ 在两者之间插值，取 $l$ 为所得带参数模空间给出的不变量，见\sr{fig:17.5}{图~17.5}。
\SourceFigure{figure-17-5}{0.51092253}{17.5}
